\section{Problem Statement}
The fundamental task in automated sleep analysis is the classification of multichannel electroencephalogram (EEG)
signals into five discrete sleep stages: Wake, N1, N2, N3, and REM.
While numerous machine learning models have been developed for this purpose, a critical and
persistent challenge is their failure to generalize to new, unseen data.

This lack of generalization arises because most high-performing models are trained and validated on limited and
homogeneous datasets, often sourced from a single clinical center with a specific patient demographic and a
consistent hardware setup. Consequently, when these models are applied to data from different populations
or recorded with different equipment their performance weakens.
This severely restricts their practical utility in real-world clinical settings, where patient
populations are inherently diverse in age, health status, and other factors that influence EEG patterns.

Therefore, the central problem addressed by this thesis is the design of a sleep staging model
that is robust to domain shift. A successful solution must not only achieve good classification accuracy
across all sleep stages, accounting for their natural class imbalance, but must also maintain this performance
consistently across datasets.
This thesis will investigate the use of the transformer architecture to learn domain-invariant
features from EEG signals to solve this generalization challenge.

